\documentclass{article}
\usepackage{amsmath,hyperref}
\begin{document}
\section*{How many staged relay sites support the distal role read?}
\paragraph{Source and scope.}
Use the full pinned Cave Circuit 01 visibility graph to compute
weighted pair costs for 20 reference sites on the distal
loop used in the whole-graph routing study. Each reference site
uses its integer SDF tile origin. Assigning a tile origin to that
tile's lookup vertex is an explicit projection model, not a proof
that a carried robot can occupy the origin, or a query of the
full historical voxel lookup. Compile the pinned range model,
visibility selection, static breadcrumb registration, broker
dispatch and receiver callback. Use the deterministic strong
signal/no-fade regime of the preceding source tests. The motion
graph's authored gate closures do not change these radio inputs.

The receiver projects to tile 103. The two possible report sources
project to 118 and 117, the reachable sides of the alternative
physical inspection cut in the separate graph model. No information
about which gate is actually closed follows merely from the source
role or radio condition. A correct gate observation would have to
originate there, and its availability and placement costs are not
established. Test one byte from each projected source, with both
qualified callbacks required for the two-case comparison.

\paragraph{Charged observation paths at matched placements.}
The direct weighted costs are 18 and 10; the Euclidean origin
distances are about 79.69 and 73.48 meters. Under the test 72-meter
range neither direct channel succeeds. Of the 17 allowed relay
origins, no one- or two-site placement serves both sources.
Six three-site placements do, including $(104,105,107)$.
At the test 80-meter range there are two successful two-site
placements; $(105,107)$ is one. The 500-meter arm has the same
finite placement counts because enlarging the range does not
remove the visibility cost of 18 for the first source. These
minima refer only to the named candidate sites and bounded sweep,
not arbitrary continuous placements or an engineered cave network.

The actual native callbacks agree on every two-source receipt mask
with the direct independent-hop feasibility calculation in this
sweep (9,642 placement/range cases). Earlier, smaller-cave cases
showed that this equality is not a source law valid for every
geometry or range. A static relay entered after the first dispatch
does not retroactively deliver the first packet; entered after both,
it cannot change either earlier observation. A moving-only name
does not register the static relay at all. All deliveries still
come from an actual application callback, not an ideal path test.

\paragraph{Deployment residuals.}
This is a source-backed check of projected observation reach, not
a paired physical-motion and radio deadline proof. Inflation of
RF range without cost modeling is not credited as an improvement.
The palette is drawn from an authored held-in loop, no early gate
sensor or actual relay travel is established, and no independent
I2 geometry is selected. Same-interface direct graph feasibility
ties in this regime.
\end{document}
